% !TeX root = ../../Book3.tex
% 纲要标题：### 9.3 有限型代数的维数
% 纲要位置：工作记录/计划纲要.md，第 2091 行。
% 纲要细目（写作范围）：
% * 多项式环的维数
% * Noether 正规化和超越次数
% * 整扩张的维数保持


\section{有限型代数的维数}
\label{b3:sec:0903}

多项式环中有明显的素理想链，但仅展示一条长链只能给出下界。上界需要证明：每增加一个真正的新方程，分式域中能保持代数独立的元素就会减少。Noether 正规化随后把一般有限型整代数与多项式环联系起来。


\subsection{多项式环的维数}
\label{b3:subsec:0903:01}

\begin{lemma}\label{b3:lem:transcendence-degree-drop}
设 \(A\) 为域 \(k\) 上的有限型整代数，\(0\ne\mathfrak p\) 为素理想。则
\[
\operatorname{trdeg}_k\operatorname{Frac}(A/\mathfrak p)
<\operatorname{trdeg}_k\operatorname{Frac}(A).
\]
\end{lemma}
\begin{proof}
从 \(A/\mathfrak p\) 的有限代数生成组中取极大代数独立组
\(\bar y_1,\ldots,\bar y_r\)，并在 \(A\) 中取提升 \(y_i\)。其他生成元在
\(k(\bar y_1,\ldots,\bar y_r)\) 上代数，故左侧超越次数是 \(r\)。提升后的 \(y_i\) 也代数独立。

取 \(0\ne a\in\mathfrak p\)。我们证明 \(a\) 在 \(k(y_1,\ldots,y_r)\) 上超越。
若它代数，取最小多项式并清除分母，得到
\[
c_s(y)a^s+\cdots+c_1(y)a+c_0(y)=0,
\qquad 0\ne c_0\in k[Y_1,\ldots,Y_r].
\]
常数项非零，因为 \(a\ne0\)，它的最小多项式不能被变量整除。模掉 \(\mathfrak p\) 得
\(c_0(\bar y)=0\)，与代数独立性矛盾。因此 \(y_1,\ldots,y_r,a\) 代数独立，右侧超越次数至少是 \(r+1\)。
\end{proof}

\begin{theorem}\label{b3:thm:polynomial-dimension}
任意域 \(k\) 及整数 \(n\ge0\) 满足
\[
\dim k[X_1,\ldots,X_n]=n.
\]
更一般地，有限型整 \(k\) 代数 \(A\) 的维数不超过
\(\operatorname{trdeg}_k\operatorname{Frac}(A)\)。
\end{theorem}
\begin{proof}
对严格素链 \(\mathfrak p_0\subsetneq\cdots\subsetneq\mathfrak p_r\)，逐次将上一引理用于
\(A/\mathfrak p_i\) 及其非零素理想 \(\mathfrak p_{i+1}/\mathfrak p_i\)。超越次数每一步至少下降一，而起点超越次数不超过 \(\operatorname{trdeg}_k\operatorname{Frac}(A)\)：商环中的代数独立元素提升后仍然代数独立。这证明上界。对于多项式环，
\[
(0)\subsetneq(X_1)\subsetneq(X_1,X_2)\subsetneq\cdots
\subsetneq(X_1,\ldots,X_n)
\]
的各项都是素理想，给出长度 \(n\) 的链。\(n=0\) 时环为域，结论同样成立。
\end{proof}

例如 \(k[X,Y,Z]/(XY)\) 的两个极小素理想为 \((\bar X)\) 和 \((\bar Y)\)，相应商环均为二元多项式环，故其维数为二。这个计算按不可约分支分别进行，不把有零因子的环擅自当成具有一个分式域的整环。

\subsection{Noether 正规化与超越次数}
\label{b3:subsec:0903:02}

\begin{theorem}\label{b3:thm:affine-dimension}\label{b3:thm:affine-dimension-trdeg}
设 \(k\) 为域。若 \(A\) 是有限型整 \(k\) 代数，则
\[
\dim A=\operatorname{trdeg}_k\operatorname{Frac}(A).
\]
若 \(A\ne0\) 只是有限型交换 \(k\) 代数，则
\[
\dim A=\max_{\mathfrak p\in\operatorname{Min}A}
\operatorname{trdeg}_k\operatorname{Frac}(A/\mathfrak p).
\]
这里 \(\operatorname{Min}A\) 表示极小素理想的集合。
\end{theorem}
\begin{proof}
整环情形由\cref{b3:thm:noether-normalization}选出代数独立元素
\(y_1,\ldots,y_d\)，使 \(A\) 是 \(B=k[y_1,\ldots,y_d]\) 上的有限模。
在 \(B\setminus\{0\}\) 处局部化后，所得环是有限维
\(k(y_1,\ldots,y_d)\) 代数且为整环，因而是域，等于 \(\operatorname{Frac}(A)\)。
所以该域在 \(k(y_1,\ldots,y_d)\) 上有限，超越次数为 \(d\)。

多项式环 \(B\) 中有长度 \(d\) 的素链。由\cref{b3:thm:lying-over}提升首项，再反复用\cref{b3:thm:going-up}提升其余各项；收缩严格增加，提升链也严格增加。因此 \(\dim A\ge d\)。反向不等式已由\cref{b3:thm:polynomial-dimension}证明。

一般情形下，\(A\) 为 Noether 环，只有有限多个极小素理想。任意素链的首项都包含某个极小素理想，因此整条链都来自某个 \(A/\mathfrak p\)。反之这些商环中的链也是 \(A\) 中的链。这给出
\(\dim A=\max_{\mathfrak p\in\operatorname{Min}A}\dim(A/\mathfrak p)\)，再应用整环情形。
\end{proof}

尖点环 \(k[t^2,t^3]\) 整包含 \(k[t^2]\)，其分式域为 \(k(t)\)，所以维数为一。
变量之间有关系 \((t^3)^2=(t^2)^3\)，两个代数生成元不意味着二维。
相反，\(k[x,y,x^{-1}]\) 虽添入了一个逆元，分式域仍为 \(k(x,y)\)，所以维数仍为二。
有限型条件在这里保证 Noether 正规化可用；不能把该公式未经说明用于任意子环。

\subsection{整扩张的维数保持}
\label{b3:subsec:0903:03}

\begin{theorem}\label{b3:thm:integral-dimension}
若交换环间的包含 \(A\subseteq B\) 是整扩张，则 \(\dim A=\dim B\)。不要求扩张有限，也不要求两环为整环。
\end{theorem}
\begin{proof}
\(B\) 中的严格素链收缩到 \(A\) 后仍严格，因为\cref{b3:thm:incomparability}说明具有相同收缩的两个可比较素理想必相等。因此 \(\dim B\le\dim A\)。
反之，从 \(A\) 中任意有限素链出发，先用 Lying-over 提升首项，再用 Going-up 逐项提升；不同的收缩保证每次包含严格。所以 \(\dim B\ge\dim A\)。这比较的是所有有限链的长度，维数无穷时也成立。
\end{proof}

该定理解释了有限映射为什么保存整体维数，却没有声称每个点的高度在任意整扩张下都保持。
若基环 \(A\) 为整闭整环，扩环 \(B\) 也为整环，则 Going-down 能补足点处的结论。
\begin{corollary}\label{b3:cor:integral-height-normal-base}
若 \(A\subseteq B\) 为整环间的整扩张，且 \(A\) 整闭，则对
\(\mathfrak q\in\operatorname{Spec}B\)、\(\mathfrak p=\mathfrak q\cap A\)，有
\(\operatorname{ht}\mathfrak q=\operatorname{ht}\mathfrak p\)。
\end{corollary}
\begin{proof}
以 \(\mathfrak q\) 为末项的链收缩后仍严格，故得到一个方向。反向由\cref{b3:thm:going-down}，从指定的 \(\mathfrak q\) 开始向下提升每条以 \(\mathfrak p\) 为末项的有限链。
\end{proof}

整同态若不是单射，应先把基环换成其像：此时能推出的是
\[
\dim B=\dim\Bigl(A/\ker(A\rightarrow B)\Bigr).
\]
例如商同态
\(k[x]\rightarrow k\)、\(x\mapsto0\) 是整同态，却并不保持原基环的维数。
